% =====================================================
% PART 3: MODEL TRAINING AND EVALUATION
% =====================================================
\chapter{Model Training and Evaluation}

\section{Random Forest principle}

A \textbf{decision tree} classifies an observation through a sequence of rules on features (e.g., ``\texttt{Connection\_Attempts} $>$ 5000?''). On its own, it tends to overfit easily.

A \textbf{Random Forest} combines 100 trees, each trained on:
\begin{itemize}[leftmargin=*,itemsep=2pt]
    \item a bootstrap sample (random sampling with replacement) of the data;
    \item a random subset of features at each split.
\end{itemize}

The final prediction is obtained through a \textbf{majority vote} among all trees, which reduces variance and overfitting compared to a single decision tree.

\section{Training}

\begin{lstlisting}[language=Python]
from sklearn.ensemble import RandomForestClassifier

model_forest = RandomForestClassifier(
    n_estimators=100, random_state=42, n_jobs=-1
)
model_forest.fit(X_train, y_train)
\end{lstlisting}

\begin{itemize}[leftmargin=*,itemsep=2pt]
    \item \texttt{n\_estimators=100}: balance between performance and computation time.
    \item \texttt{random\_state=42}: reproducibility.
    \item \texttt{n\_jobs=-1}: uses all CPU cores for parallel training.
\end{itemize}

\section{Test set evaluation}

\begin{lstlisting}[language=Python]
from sklearn.metrics import accuracy_score, confusion_matrix, classification_report

y_pred = model_forest.predict(X_test)
print(accuracy_score(y_test, y_pred))
print(confusion_matrix(y_test, y_pred))
print(classification_report(y_test, y_pred))
\end{lstlisting}

\begin{table}[H]
\centering
\caption{Confusion Matrix (test set, 30,900 samples)}
\label{tab:confusion}
\begin{tabular}{cc>{\centering\arraybackslash}p{3cm}>{\centering\arraybackslash}p{3cm}}
\toprule
 & & \multicolumn{2}{c}{\textbf{Predicted}} \\
 & & Normal (0) & Attack (1) \\
\midrule
\multirow{2}{*}{\textbf{Actual}} & Normal (0) & 19 864 & 2 \\
 & Attack (1) & 0 & 11 034 \\
\bottomrule
\end{tabular}
\end{table}

\begin{table}[H]
\centering
\caption{Classification metrics}
\label{tab:report}
\begin{tabular}{lcccc}
\toprule
\textbf{Class} & \textbf{Precision} & \textbf{Recall} & \textbf{F1-score} & \textbf{Support} \\
\midrule
0 (Normal) & 1.00 & 1.00 & 1.00 & 19 866 \\
1 (Attack) & 1.00 & 1.00 & 1.00 & 11 034 \\
\midrule
\textbf{Accuracy} & \multicolumn{4}{c}{99.99\% (30,900 samples)} \\
\bottomrule
\end{tabular}
\end{table}

\textbf{Result:} accuracy of 99.99\%, 2 false positives, 0 false negatives. Precision/recall/F1-score = 1.00 for both classes.

\subsection*{Discussion}

Such a near-perfect score should be interpreted with caution: it may reflect a very clear separation between classes in this dataset (low noise, limited overlap) rather than guaranteed real-world performance. No data leakage was detected (IP addresses excluded from $X$), but this result justifies validation on a second dataset (Part 4).

\section{Confusion matrix visualization}

\begin{lstlisting}[language=Python]
import matplotlib.pyplot as plt
import seaborn as sns
import numpy as np

fig, axes = plt.subplots(1, 2, figsize=(15, 6))

cm = confusion_matrix(y_test, y_pred)
sns.heatmap(cm, annot=True, fmt='d', cmap='Blues', ax=axes[0],
            xticklabels=['Normal', 'Attack'], yticklabels=['Normal', 'Attack'])
axes[0].set_title("Confusion Matrix")

plt.tight_layout()
plt.savefig("figures/confusion_et_importance.png", dpi=150)
plt.show()
\end{lstlisting}

\figplaceholder{confusion_et_importance.png}{Confusion matrix and top 10 most important features}{0.95}

Feature importance analysis helps identify the most discriminative network indicators (e.g., connection attempts, TCP flags, throughput) and confirms that the model relies on meaningful behavioral signals rather than artifacts.
